%% kapa: exam=1 round=37 year=2026 period=2 form=A subject=accounting subject_ko=회계학 seq=36 number=76
%% source: Q-Net 감정평가사 기출문제 (cst00309) artlSeq=5258924 file=raw/1008/5258924/2026년 제37회 감정평가사 1차 2교시.pdf
%% conversion: kordoc 4.24.51 -> tools/gen1.py (auto); review=proofread-sonnet; flags=none
\begin{question}{76}
종합원가계산에 관한 설명으로 옳지 않은 것은?
\begin{choices}
  \item 작업별로 원가계산이 이루어지는 개별원가계산과 달리 종합원가계산에서는 각 제조공정별로 제조원가를 집계하고, 일정기간동안 해당 공정에서 생산된 완성품과 기말재공품에 제조원가를 배분하여 당기제품제조원가와 기말재공품원가를 계산한다.
  \item 기초재공품이 존재하는 경우에는 재공품계정의 차변에 집계된 제조원가가 전기의 제조원가와 당기에 발생한 제조원가로 구성되므로, 이를 완성품과 기말재공품에 배분하기 위해서는 원가흐름의 가정이 필요하다.
  \item 공손품이 처분가치를 갖는 경우에는 공손품의 순실현가능가치를 공손품이라는 자산계정으로 기록하고, 정상공손원가에서 정상공손의 순실현가능가치를 차감한 순정상공손원가만을 합격품에 배분한다.
  \item 기말재공품이 당기 중에 검사시점을 통과한 경우, 정상공손의 원가는 완성품과 기말재공품의 물량을 기준으로 완성품과 기말재공품에 배분한다.
  \item 정상공손과 비정상공손을 물리적으로 구분하는 것은 곤란하지만, 각 원가요소별 정상공손과 비정상공손의 완성품환산량 단위당 원가가 상이하기 때문에 정상공손의 원가는 합격품에 배분하는 반면, 비정상공손의 원가는 당기 기간비용으로 처리한다.
\end{choices}
\end{question}
\answer{5}
